\chapter{Metrics-Driven Acceptance Criteria}
\label{ch:metrics-driven}

Acceptance criteria grounded in metrics ensure teams ship outcomes, not just outputs. This chapter explains how to choose metrics, balance trade-offs, and operationalize measurement without inciting metric gaming.

\section{The Case for Metrics-Driven Acceptance}
\label{sec:case-metrics}

Metrics remove ambiguity and align stakeholders on what success means.

\subsection{From Subjective to Objective}
Instead of debating whether ``the experience feels better,'' teams can agree to ``activation rate increases by 8\% with p \textless{} 0.05.'' Data replaces opinion, expediting decisions.

\subsection{Enabling Automation}
When criteria map to metrics, CI/CD pipelines can gate deployments, and dashboards can alert teams when performance drifts. Automation shortens feedback loops.

\subsection{Learning and Improvement}
Tracking acceptance metrics over time reveals whether improvements persist, enabling continuous optimization and reducing regression risk.

\section{Choosing Metrics}
\label{sec:choosing-metrics}

Selecting the right metrics is critical.

\subsection{Desirable Properties of Metrics}
Metrics should be measurable, relevant, actionable, timely, and controllable. Avoid metrics that rely on data you cannot reliably capture.

\subsection{Leading vs. Lagging Indicators}
Leading indicators predict outcome trends (e.g., onboarding completions), while lagging indicators measure final impact (e.g., revenue). Balance both to catch issues early without losing sight of ultimate goals.

\subsection{Input, Output, and Outcome Metrics}
Inputs measure effort (tests written), outputs capture deliverables (features shipped), and outcomes measure impact (retention). Prioritize outcomes but use input/output metrics when necessary to proxy future impact.

\subsection{Metric Selection Framework}
Ask: What behavior do we want to influence? What decision will the metric inform? Who owns it? Piloting metrics on a small cohort before formal adoption reduces surprises.

\section{Metrics for Feature Development}
\label{sec:feature-metrics}

Features require a mix of user and technical metrics.

\subsection{Usage Metrics}
Adoption rate, active users, or funnel completion quantifies whether users try the feature. Define thresholds such as ``50\% of eligible users interact within 30 days.''

\subsection{Engagement Metrics}
Depth metrics (time on task, frequency) ensure users derive value, not just trial usage.

\subsection{Satisfaction Metrics}
CSAT, NPS, or qualitative survey responses confirm that changes align with expectations. Combine numeric targets with qualitative feedback quotas.

\subsection{Performance Metrics}
Latency, error rate, and availability guard against regressions. Acceptance might require ``P95 latency $\le 200\,\text{ms}$ with error rate $<0.1\%$.''

\subsection{Example: New Dashboard Feature}
Acceptance criteria could include: 40\% adoption among analysts, +10 point satisfaction lift, P95 load $\le 1.5$ seconds, and no increase in support tickets.

\section{Metrics for ML Models}
\label{sec:ml-metrics}

Model launches need offline, online, and business metrics.

\subsection{Offline Evaluation Metrics}
Specify thresholds for precision, recall, F1, AUC, or regression metrics. Include confidence intervals and fairness constraints.

\subsection{Online Evaluation Metrics}
A/B tests validate live impact—click-through, conversion, or churn. Acceptance requires statistically significant lift without violating guardrails.

\subsection{Model Health Metrics}
Monitor latency, resource usage, feature drift, and calibration. Include acceptable ranges and alert thresholds in requirements.

\subsection{Business Impact Metrics}
Tie model performance to revenue, cost savings, or risk reduction. ``Reduce fraud losses by \$500K/month'' resonates with executives.

\subsection{Example: Recommendation System}
Criteria might include offline NDCG@10 $\ge 0.48$, online CTR $\ge 4\%$, diversity index $\ge 0.3$, latency $\le 100\,\text{ms}$, and no increase in churn.

\section{Metrics for Experiments}
\label{sec:experiment-metrics}

Experiments need clear success metrics and guardrails.

\subsection{Primary, Secondary, and Guardrail Metrics}
Primary metrics determine success; secondary metrics provide context; guardrails ensure you do not harm adjacent KPIs. Document all before launch.

\subsection{Statistical Power and Duration}
Acceptance should specify minimum detectable effect, required sample size, and planned duration to avoid underpowered tests.

\subsection{Decision Criteria}
Define thresholds and escalation paths: ``If the 95\% CI includes zero, rerun or pivot.''

\section{Metrics for Infrastructure and Platform Work}
\label{sec:infra-metrics}

Non-user-facing work still needs measurable outcomes.

\subsection{Reliability Metrics}
Use SLOs/SLA (availability, latency) and error budgets as acceptance criteria.

\subsection{Efficiency Metrics}
Measure resource utilization, build times, or deployment frequency improvements.

\subsection{Developer Experience Metrics}
Track satisfaction surveys, support tickets, or onboarding time for internal platforms.

\section{Balancing Multiple Metrics}
\label{sec:balancing-metrics}

Most work involves trade-offs.

\subsection{Identifying Trade-offs}
Document conflicting goals (precision vs. recall, latency vs. accuracy) and quantify acceptable ranges.

\subsection{Weighted Scoring}
Create weighted composite scores when decisions require multiple metrics. For example, assign 60\% weight to relevance, 25\% to diversity, 15\% to latency.

\subsection{Pareto Optimality}
Visualize Pareto frontiers to compare solutions without forcing a single scalar metric.

\subsection{Decision Frameworks}
Facilitate discussions across stakeholders about how to react when metrics conflict. Document escalation paths.

\subsection{Example: Search Ranking Model}
List acceptance thresholds for relevance, diversity, personalization, and latency. Provide weighting to choose among candidate models.

\section{Avoiding Metric Pitfalls}
\label{sec:metric-pitfalls}

Metrics are powerful but can be misused.

\subsection{Goodhart's Law}
When metrics become targets, behavior can distort them. Use guardrails and rotate metrics periodically to reduce gaming.

\subsection{Vanity Metrics}
Avoid metrics that look good but lack business value (e.g., total sign-ups without engagement). Prioritize actionable metrics tied to outcomes.

\subsection{Local vs. Global Optimization}
Optimize metrics holistically. Example: focusing solely on conversion might increase refunds; include profitability guardrails.

\subsection{Overfitting to Metrics}
Do not tailor solutions narrowly to measured metrics; maintain qualitative reviews and fail-safes.

\subsection{Ignoring Context and Nuance}
Metrics cannot capture all nuance. Document how qualitative insights or ethics considerations can override metrics when necessary.

\section{Metrics Dashboards and Monitoring}
\label{sec:dashboards}

Visibility keeps metrics actionable.

\subsection{Dashboard Design Principles}
Dashboards should prioritize clarity, highlight deviations, and offer drill-downs. Tailor detail to the audience.

\subsection{Real-Time Monitoring}
Implement alerting when metrics breach thresholds. Automate rollback triggers where appropriate.

\subsection{Historical Trends}
Track seasonal patterns, moving averages, and regression lines to contextualize changes.

\subsection{Automated Analysis}
Leverage anomaly detection or forecasting to flag issues early. Record annotations when changes occur for future audits.

\subsection{Example: Model Performance Dashboard}
Provide a layout showing offline metrics, online experiment outcomes, drift detectors, and business KPIs, all linked to relevant tickets.

\section{Metrics in Continuous Improvement}
\label{sec:continuous-improvement}

Metrics underpin learning organizations.

\subsection{Baseline, Measure, Improve, Repeat}
Document baselines before changes, measure results, and iterate. Incorporate metrics into retrospectives to assess whether processes improved.

\subsection{Improvement Targets}
Set realistic stretch goals. Use statistical process control charts to track whether improvements are sustained.

\subsection{Experimentation Culture}
Make experimentation the default: small, instrumented changes inform future decisions. Encourage sharing of negative results.

\subsection{Sharing and Learning}
Create forums where teams present metric outcomes, case studies, and lessons learned. Institutionalize best practices via documentation and training.

\section{Summary}
\label{sec:metrics-summary}

Metrics-driven acceptance criteria bring clarity, accountability, and learning to the product lifecycle. By choosing meaningful metrics, balancing trade-offs, and monitoring outcomes, teams ensure that shipped features deliver tangible value. The next chapter discusses documentation patterns that preserve these insights.